\Article{22}
We shall obtain the variation of the whole surface $U$ by integrating this expression over all the elements $\mathrm{d}U$. For this purpose we shall treat the two parts of this integral separately, namely
\[
\int\mathrm{d}U.((\eta\eta+\zeta\zeta)\frac{\mathrm{d}\delta x}{\mathrm{d}x}-\xi\eta\frac{\mathrm{d}\delta y}{\mathrm{d}x}-\xi\zeta\frac{\mathrm{d}\delta z}{\mathrm{d}x})=A
\]
and
\[
\int\mathrm{d}U(-\xi\eta\frac{\mathrm{d}\delta x}{\mathrm{d}y}+(\xi\xi+\zeta\zeta)\frac{\mathrm{d}\delta y}{\mathrm{d}y}-\eta\zeta\frac{\mathrm{d}\delta z}{\mathrm{d}y})=B
\]

Imagine a plane normal to the coordinate axis $y$, chosen so that its fixed value of $y$ lies within the range bounded by the extreme values of $y$ on the surface $U$. This plane will cut the boundary $P$ in two, four, six, and so on, points. Their first coordinates in order are $x^0,x',x''$ etc. Let the other quantities belonging to these points be distinguished by corresponding indices. In the same way let the surface be cut by another plane, infinitely close and parallel to the first, with second coordinate $y+\mathrm{d}y$. Between these planes lie the boundary elements $\mathrm{d}P^0$, $\mathrm{d}P'$, $\mathrm{d}P''$ etc., and it is easy to see that
\[
\mathrm{d}y=-Y^0\mathrm{d}P^0=+Y'\mathrm{d}P'=-Y''\mathrm{d}P''=+Y'''\mathrm{d}P'''\quad\text{etc.}
\]
\SourceLeaf{65}If we further imagine infinitely many planes normal to the coordinate axis $x$, each element $\mathrm{d}x$ between $x^0$ and $x'$, or between $x''$ and $x'''$ etc., will have a corresponding element $\mathrm{d}U=\frac{\mathrm{d}x.\mathrm{d}y}{\zeta}$. It is thus clear that the part of the integral $A$ corresponding to the part of the surface between the planes $y,y+\mathrm{d}y$ is obtained by the integration
\[
\mathrm{d}y\int\mathrm{d}x.(\frac{\eta\eta+\zeta\zeta}{\zeta}.\frac{\mathrm{d}\delta x}{\mathrm{d}x}-\frac{\xi\eta}{\zeta}.\frac{\mathrm{d}\delta y}{\mathrm{d}x}-\xi\frac{\mathrm{d}\delta z}{\mathrm{d}x})
\]
extended from $x=x^0$ to $x=x'$, then from $x=x''$ to $x=x'''$ etc. The indefinite integral is expressed by
\[
\begin{aligned}
&(\frac{\eta\eta+\zeta\zeta}{\zeta}.\delta x-\frac{\xi\eta}{\zeta}.\delta y-\xi\delta z)\mathrm{d}y\\
&-\mathrm{d}y\int(\delta x.\frac{\mathrm{d}\frac{\eta\eta+\zeta\zeta}{\zeta}}{\mathrm{d}x}-\delta y.\frac{\mathrm{d}\frac{\xi\eta}{\zeta}}{\mathrm{d}x}-\delta z.\frac{\mathrm{d}\xi}{\mathrm{d}x})\mathrm{d}x
\end{aligned}
\]
from which in our case we obtain
\[
\begin{aligned}
&(\frac{\eta^0\eta^0+\zeta^0\zeta^0}{\zeta^0}.\delta x^0-\frac{\xi^0\eta^0}{\zeta^0}.\delta y^0-\xi^0\delta z^0)Y^0\mathrm{d}P^0\\
+&(\frac{\eta'\eta'+\zeta'\zeta'}{\zeta'}.\delta x'-\frac{\xi'\eta'}{\zeta'}.\delta y'-\xi'\delta z')Y'\mathrm{d}P'\\
+&(\frac{\eta''\eta''+\zeta''\zeta''}{\zeta''}.\delta x''-\frac{\xi''\eta''}{\zeta''}.\delta y''-\xi''\delta z'')Y''\mathrm{d}P''\\
+&\text{etc.}\\
-&\int\zeta\mathrm{d}U.(\delta x.\frac{\mathrm{d}\frac{\eta\eta+\zeta\zeta}{\zeta}}{\mathrm{d}x}-\delta y.\frac{\mathrm{d}\frac{\xi\eta}{\zeta}}{\mathrm{d}x}-\delta z.\frac{\mathrm{d}\xi}{\mathrm{d}x})
\end{aligned}
\]
or, equivalently,
\[
\begin{aligned}
&\Sigma(\frac{\eta\eta+\zeta\zeta}{\zeta}.\delta x-\frac{\xi\eta}{\zeta}.\delta y-\xi\delta z)Y\mathrm{d}P\\
&-\int\zeta\mathrm{d}U.(\delta x.\frac{\mathrm{d}\frac{\eta\eta+\zeta\zeta}{\zeta}}{\mathrm{d}x}-\delta y.\frac{\mathrm{d}\frac{\xi\eta}{\zeta}}{\mathrm{d}x}-\delta z.\frac{\mathrm{d}\xi}{\mathrm{d}x})
\end{aligned}
\]
where the summation must extend over all the elements $\mathrm{d}P$ and the integration over all the elements $\mathrm{d}U$ between the planes $y$ and $y+\mathrm{d}y$.

The whole quantity $A$ is thus expressed by
\[
\begin{aligned}
&\int(\frac{\eta\eta+\zeta\zeta}{\zeta}.\delta x-\frac{\xi\eta}{\zeta}.\delta y-\xi\delta z)Y\mathrm{d}P\\
&-\int\zeta\mathrm{d}U(\delta x.\frac{\mathrm{d}\frac{\eta\eta+\zeta\zeta}{\zeta}}{\mathrm{d}x}-\delta y.\frac{\mathrm{d}\frac{\xi\eta}{\zeta}}{\mathrm{d}x}-\delta z.\frac{\mathrm{d}\xi}{\mathrm{d}x})
\end{aligned}
\]
where the first integration must extend over the whole boundary $P$ and the second over the whole surface $U$.

\SourceLeaf{66}\Article{23}
By exactly analogous reasoning we find
\[
\begin{aligned}
B={}&\int(\frac{\xi\eta}{\zeta}.\delta x-\frac{\xi\xi+\zeta\zeta}{\zeta}.\delta y+\eta\delta z)X\mathrm{d}P\\
&+\int\zeta\mathrm{d}U(\delta x.\frac{\mathrm{d}\frac{\xi\eta}{\zeta}}{\mathrm{d}y}-\delta y.\frac{\mathrm{d}\frac{\xi\xi+\zeta\zeta}{\zeta}}{\mathrm{d}y}+\delta z.\frac{\mathrm{d}\eta}{\mathrm{d}y})
\end{aligned}
\]
Thus, setting for any point of the boundary $P$
\[
\begin{aligned}
&[X\xi\eta+Y(\eta\eta+\zeta\zeta)]\delta x-[X(\xi\xi+\zeta\zeta)+Y\xi\eta]\delta y\\
&+(X\eta\zeta-Y\xi\zeta)\delta z=\zeta Q
\end{aligned}
\]
and for any point of the surface $U$
\[
\begin{aligned}
&(\frac{\mathrm{d}\frac{\xi\eta}{\zeta}}{\mathrm{d}y}-\frac{\mathrm{d}\frac{\eta\eta+\zeta\zeta}{\zeta}}{\mathrm{d}x})\zeta\delta x+(\frac{\mathrm{d}\frac{\xi\eta}{\zeta}}{\mathrm{d}x}-\frac{\mathrm{d}\frac{\xi\xi+\zeta\zeta}{\zeta}}{\mathrm{d}y})\zeta\delta y\\
&+(\frac{\mathrm{d}\xi}{\mathrm{d}x}+\frac{\mathrm{d}\eta}{\mathrm{d}y})\zeta\delta z=V
\end{aligned}
\]
we finally have
\[
\delta U=\int Q\mathrm{d}P+\int V\mathrm{d}U
\]
where the first integration must extend over the whole boundary $P$ and the second over the whole surface $U$.

\Article{24}
The formulas just given for $Q$ and $V$ can be shortened considerably. With the help of the equation $X\xi+Y\eta+Z\zeta=0$, $Q$ at once takes the following symmetric form.
\[
Q=(Y\zeta-Z\eta)\delta x+(Z\xi-X\zeta)\delta y+(X\eta-Y\xi)\delta z
\]

To put the expression obtained for $V$ into a neater form as well, we observe that from the formulas
\[
\frac{\mathrm{d}z}{\mathrm{d}x}=-\frac{\xi}{\zeta},\qquad \frac{\mathrm{d}z}{\mathrm{d}y}=-\frac{\eta}{\zeta}
\]
it follows that
\[
\frac{\mathrm{d}\frac{\xi}{\zeta}}{\mathrm{d}y}=\frac{\mathrm{d}\frac{\eta}{\zeta}}{\mathrm{d}x}
\]
Hence
\[
\frac{\mathrm{d}\frac{\xi\eta}{\zeta}}{\mathrm{d}y}=\frac{\xi}{\zeta}.\frac{\mathrm{d}\eta}{\mathrm{d}y}+\eta\frac{\mathrm{d}\frac{\xi}{\zeta}}{\mathrm{d}y}=\frac{\xi}{\zeta}.\frac{\mathrm{d}\eta}{\mathrm{d}y}+\eta\frac{\mathrm{d}\frac{\eta}{\zeta}}{\mathrm{d}x}
\]
\SourceLeaf{67}Further, from $\xi\xi+\eta\eta+\zeta\zeta=1$ we deduce
\[
\xi\frac{\mathrm{d}\xi}{\mathrm{d}x}+\eta\frac{\mathrm{d}\eta}{\mathrm{d}x}+\zeta\frac{\mathrm{d}\zeta}{\mathrm{d}x}=0
\]
and from this
\[
\frac{\mathrm{d}\frac{\eta\eta+\zeta\zeta}{\zeta}}{\mathrm{d}x}=\eta\frac{\mathrm{d}\frac{\eta}{\zeta}}{\mathrm{d}x}+\frac{\eta}{\zeta}.\frac{\mathrm{d}\eta}{\mathrm{d}x}+\frac{\mathrm{d}\zeta}{\mathrm{d}x}=\eta\frac{\mathrm{d}\frac{\eta}{\zeta}}{\mathrm{d}x}-\frac{\xi}{\zeta}.\frac{\mathrm{d}\xi}{\mathrm{d}x}
\]
Substituting these values in the coefficient of $\delta x$ in the expression for $V$, we find that it becomes
\[
=\xi(\frac{\mathrm{d}\xi}{\mathrm{d}x}+\frac{\mathrm{d}\eta}{\mathrm{d}y})
\]
In exactly the same way the coefficient of $\delta y$ in the same expression becomes
\[
\eta(\frac{\mathrm{d}\xi}{\mathrm{d}x}+\frac{\mathrm{d}\eta}{\mathrm{d}y})
\]
We thus obtain
\[
V=(\xi\delta x+\eta\delta y+\zeta\delta z)(\frac{\mathrm{d}\xi}{\mathrm{d}x}+\frac{\mathrm{d}\eta}{\mathrm{d}y})
\]

